% why/vision.tex

\chapter{Vision}

% why/sections/why_railway_is_different.tex

\section*{Why Railway is Different}

\begin{insight}
Designing a railway alignment is fundamentally different from designing
a road.

A road vehicle can steer.

A train cannot.

\medskip

A road alignment may tolerate small geometric imperfections.
A railway alignment cannot.

\medskip

For a train:
\begin{itemize}
\item curvature directly determines lateral acceleration,
\item curvature variation determines jerk,
\item and both directly affect safety, comfort, and wear.
\end{itemize}

\medskip

As a consequence:
\begin{itemize}
\item geometry is not just shape,
\item curvature is not just a derived quantity,
\item and transitions are not optional smoothing elements.
\end{itemize}

\medskip

They are the system.
\end{insight}


\section{From Problem to Paradigm}

\begin{remark}
The previous chapters have identified a fundamental limitation of current
alignment modelling approaches: the fragmentation of geometry, dynamics,
and data.
\end{remark}

\begin{remark}
This motivates not only incremental improvements, but a shift in
perspective.
\end{remark}

\section{Core Idea}

\begin{remark}
The central idea of this work is to treat railway alignment as a
\emph{state trajectory} governed by curvature and its associated
quantities.
\end{remark}

\begin{remark}
Instead of describing alignment as a sequence of geometric elements,
we propose to describe it as a function-based object:
\[
(\kappa(s), \phi(s)).
\]
\end{remark}

\begin{remark}
In this view:
\begin{itemize}
\item curvature defines geometry,
\item cant defines dynamic balancing,
\item and their interaction determines physical behaviour.
\end{itemize}
\end{remark}

\section{Alignment as a Controlled System}

\begin{remark}
The alignment is interpreted as a controlled system evolving along
arc length:
\[
\begin{aligned}
\gamma'(s) &= (\cos\theta(s), \sin\theta(s)), \\
\theta'(s) &= \kappa(s), \\
\phi'(s)   &= \omega(s).
\end{aligned}
\]
\end{remark}

\begin{remark}
Here:
\begin{itemize}
\item \(\kappa(s)\) acts as a geometric control,
\item \(\omega(s)\) governs cant evolution,
\item and the resulting state determines both geometry and dynamics.
\end{itemize}
\end{remark}

\begin{remark}
This formulation establishes a direct link between alignment design and
optimal control theory.
\end{remark}

\section{Transitions as Primary Objects}

\begin{remark}
In classical models, transitions are treated as auxiliary elements
connecting geometric primitives.
\end{remark}

\begin{remark}
In the proposed framework, transitions become the primary modelling
objects, as they define the evolution of curvature.
\end{remark}

\begin{remark}
This shifts the focus:
\begin{itemize}
\item from static geometry
\item to controlled change of state.
\end{itemize}
\end{remark}

\section{Integration of Reality}

\begin{remark}
A central requirement of the proposed approach is compatibility with
real-world data.
\end{remark}

\begin{remark}
This includes:
\begin{itemize}
\item measurement data with uncertainty,
\item heterogeneous coordinate systems (CRS),
\item legacy alignment descriptions,
\item and complex track configurations.
\end{itemize}
\end{remark}

\begin{remark}
The model must therefore support:
\begin{itemize}
\item robust reconstruction from imperfect data,
\item explicit handling of coordinate transformations,
\item and consistent representation across scales.
\end{itemize}
\end{remark}

\section{Unified Optimization Framework}

\begin{remark}
Within this paradigm, alignment design is formulated as an optimization
problem:
\[
\min J(\kappa,\phi)
\]
subject to:
\begin{itemize}
\item state dynamics,
\item engineering constraints,
\item and boundary conditions.
\end{itemize}
\end{remark}

\begin{remark}
This allows:
\begin{itemize}
\item direct incorporation of comfort criteria,
\item integration of dynamic effects,
\item and systematic exploration of design alternatives.
\end{itemize}
\end{remark}

\section{Towards Alignment-Based Information Modelling}

\begin{remark}
The proposed framework forms the basis of
\emph{alignment-based information modelling} (AIM).
\end{remark}

\begin{remark}
AIM is characterized by:
\begin{itemize}
\item curvature-driven representation,
\item state-based modelling,
\item integration of geometry and dynamics,
\item and compatibility with data-driven workflows.
\end{itemize}
\end{remark}

\begin{remark}
In this sense, alignment is not a static object, but a structured,
computable entity.
\end{remark}

\section{Implications}

\begin{remark}
Adopting this perspective has several implications:
\begin{itemize}
\item design becomes a problem of state control,
\item transitions emerge from functional requirements,
\item optimization replaces iterative manual adjustment,
\item and modelling becomes directly compatible with computation.
\end{itemize}
\end{remark}

\section{Long-Term Perspective}

\begin{remark}
The proposed framework is intended as a foundation for future
developments, including:
\begin{itemize}
\item integration with BIM systems,
\item extension to network-level modelling,
\item incorporation of vehicle dynamics,
\item and real-time optimization.
\end{itemize}
\end{remark}

\begin{remark}
It is not a finished solution, but a structured starting point.
\end{remark}

\section{Summary}

\begin{summary}
This work proposes a paradigm shift from element-based alignment
modelling to a curvature-driven, state-based representation.

By treating alignment as a controlled system and integrating geometry,
dynamics, and data, it establishes a unified foundation for modelling
and optimization.

This perspective forms the basis of alignment-based information
modelling (AIM).
\end{summary}
